% !TEX program = xelatex
\documentclass[10pt,letterpaper,twocolumn]{article}

% === GEOMETRY ===
\usepackage[top=0.75in, bottom=0.75in, left=0.75in, right=0.75in]{geometry}

% === FONTS ===
\usepackage{fontspec}
\usepackage{unicode-math}

\setmainfont{Latin Modern Roman}
\setsansfont{Latin Modern Sans}

% Custom AC fonts
\newfontfamily\acbold{ywft-processing-bold}[
  Path=../../system/public/type/webfonts/,
  Extension=.ttf
]
\newfontfamily\aclight{ywft-processing-light}[
  Path=../../system/public/type/webfonts/,
  Extension=.ttf
]
\setmonofont{Latin Modern Mono}[Scale=0.85]

% === PACKAGES ===
\usepackage{xcolor}
\usepackage{titlesec}
\usepackage{enumitem}
\usepackage{booktabs}
\usepackage{tabularx}
\usepackage{multicol}
\usepackage{fancyhdr}
\usepackage{hyperref}
\usepackage{graphicx}
\usepackage{ragged2e}
\usepackage{microtype}
\usepackage{listings}
\usepackage{natbib}

% === COLORS (AC palette) ===
\definecolor{acpink}{RGB}{180,72,135}
\definecolor{acpurple}{RGB}{120,80,180}
\definecolor{acdark}{RGB}{64,56,74}
\definecolor{acgray}{RGB}{119,119,119}

% === HYPERREF ===
\hypersetup{
  colorlinks=true,
  linkcolor=acpurple,
  urlcolor=acpurple,
  citecolor=acpurple,
  pdfauthor={Jeffrey Alan Scudder},
  pdftitle={KidLisp: A Minimal Lisp for Generative Art},
}

% === SECTION FORMATTING ===
\titleformat{\section}
  {\normalfont\bfseries\normalsize\uppercase}
  {\thesection.}
  {0.5em}
  {}
\titlespacing{\section}{0pt}{1.2em}{0.3em}

\titleformat{\subsection}
  {\normalfont\bfseries\small}
  {\thesubsection}
  {0.5em}
  {}
\titlespacing{\subsection}{0pt}{0.8em}{0.2em}

% === HEADER/FOOTER ===
\pagestyle{fancy}
\fancyhf{}
\renewcommand{\headrulewidth}{0pt}
\fancyfoot[C]{\footnotesize\thepage}

% === CUSTOM COMMANDS ===
\newcommand{\acdot}{{\color{acpink}.}}
\newcommand{\kidlisp}{\textsc{KidLisp}}

% === LISTINGS ===
\lstdefinelanguage{kidlisp}{
  morekeywords={wipe,ink,line,box,circle,tri,plot,flood,shape,zoom,scroll,spin,blur,contrast,def,let,if,once,repeat,later,embed,layer,width,height,frame,time,random,sin,cos,wiggle,melody,mic},
  sensitive=true,
  morecomment=[l]{;},
  morestring=[b]",
  literate={(}{{{\color{acgray}(}}}{1}
           {)}{{{\color{acgray})}}}{1},
}
\lstset{
  language=kidlisp,
  basicstyle=\ttfamily\small,
  keywordstyle=\color{acpurple}\bfseries,
  commentstyle=\color{acgray}\itshape,
  stringstyle=\color{acpink},
  breaklines=true,
  frame=single,
  rulecolor=\color{acgray!30},
  backgroundcolor=\color{acgray!5},
  xleftmargin=0.5em,
  xrightmargin=0.5em,
  aboveskip=0.5em,
  belowskip=0.5em,
}

% === LIST SETTINGS ===
\setlist[itemize]{nosep, leftmargin=1.2em, itemsep=0.1em}
\setlist[enumerate]{nosep, leftmargin=1.2em}

% === COLUMN SEPARATION ===
\setlength{\columnsep}{1.8em}

% === PARAGRAPH SETTINGS ===
\setlength{\parindent}{1em}
\setlength{\parskip}{0.3em}

\begin{document}

% ============ TITLE BLOCK ============

\twocolumn[{%
\begin{center}
{\acbold\fontsize{24pt}{28pt}\selectfont\color{acdark} KidLisp}\par
\vspace{0.2em}
{\aclight\fontsize{11pt}{13pt}\selectfont\color{acpink} A Minimal Lisp for Generative Art}\par
\vspace{0.6em}
{\normalsize Jeffrey Alan Scudder}\par
{\small\color{acgray} Independent Researcher / Aesthetic Computer}\par
{\small\color{acgray} ORCID: \href{https://orcid.org/0009-0007-4460-4913}{0009-0007-4460-4913}}\par
\vspace{0.3em}
{\small\color{acpurple} \url{https://kidlisp.com} \enspace$\cdot$\enspace \url{https://aesthetic.computer}}\par
\vspace{0.6em}
\rule{\textwidth}{1.5pt}
\vspace{0.5em}
\end{center}

\begin{quote}
\small\noindent\textbf{Abstract.}
\kidlisp{} is a minimal Lisp dialect with 118 built-in functions designed for creating generative art and interactive experiences. It runs inside the Aesthetic Computer platform, a mobile-first runtime for creative computing. The language prioritizes immediate visual output, minimal syntax, and social distribution over computational generality. Programs are shareable by three-character codes, embeddable in other programs, and executable at a URL. We describe the language design, its evaluator architecture, and report on adoption: 59 active authors have created 16,174 programs, forming a naturalistic corpus of creative coding practices. We argue that \kidlisp{} demonstrates a viable design point for domain-specific languages that trade expressive power for accessibility and social participation.
\end{quote}
\vspace{0.5em}
}]

% ============ 1. INTRODUCTION ============

\section{Introduction}

Creative coding environments typically require users to learn a general-purpose programming language before producing visual or sonic output. Processing~\citep{reas2007processing} uses Java syntax; p5.js~\citep{mccarthy2015p5js} uses JavaScript; Sonic Pi~\citep{aaron2016sonic} uses Ruby-like constructs. While these tools have dramatically lowered the barrier to computational expression, they still demand familiarity with variables, functions, control flow, and---in many cases---an IDE or local development environment.

\kidlisp{} takes a different approach. Rather than simplifying a general-purpose language, it starts from the question: \emph{what is the smallest language that can produce interesting generative art?} The answer is a Lisp with 118 carefully chosen built-in functions, no user-defined functions beyond \texttt{def}, no file I/O or networking, and a browser-native runtime that requires no installation.

The design is motivated by observations from No Paint~(2020), a pixel art tool whose community of non-technical users demonstrated that people learn computing through social participation in software they love. \kidlisp{} extends this insight into a programming language: every program is immediately shareable by a three-character code, viewable at a URL, and embeddable inside other programs.

This paper describes the language design~(\S\ref{sec:design}), evaluator architecture~(\S\ref{sec:architecture}), the social distribution system~(\S\ref{sec:social}), and adoption data~(\S\ref{sec:evaluation}).

% ============ 2. LANGUAGE DESIGN ============

\section{Language Design}
\label{sec:design}

\kidlisp{} is an S-expression language. Every program is a sequence of expressions evaluated top-to-bottom, producing side effects on an HTML Canvas 2D context, a Web Audio graph, or both.

\subsection{Design Principles}

\begin{enumerate}
  \item \textbf{Every expression draws.} Valid \kidlisp{} produces visible output. There is no \texttt{main}, no imports, no boilerplate.
  \item \textbf{Flat over deep.} Programs are composed by embedding other programs, not by defining abstractions. This keeps code readable at a glance.
  \item \textbf{Safe by construction.} No file I/O, no networking, no mutation of external state. Programs are sandboxed by the language's lack of dangerous primitives.
  \item \textbf{Social by default.} Every program gets a short code and a URL. Sharing is not an afterthought; it is the distribution mechanism.
\end{enumerate}

\subsection{Function Categories}

The 118 built-in functions span 12 categories:

\begin{table}[h]
\small
\centering
\begin{tabularx}{\columnwidth}{lrX}
\toprule
\textbf{Category} & \textbf{n} & \textbf{Examples} \\
\midrule
Graphics & 9 & \texttt{wipe}, \texttt{ink}, \texttt{line}, \texttt{box}, \texttt{circle} \\
Transformations & 11 & \texttt{scroll}, \texttt{zoom}, \texttt{spin}, \texttt{blur}, \texttt{suck} \\
Math & 14 & \texttt{+}, \texttt{sin}, \texttt{cos}, \texttt{random}, \texttt{wiggle} \\
Colors & 19 & CSS names + \texttt{rainbow}, \texttt{zebra} \\
System & 9 & \texttt{width}, \texttt{height}, \texttt{frame}, \texttt{fps} \\
3D & 8 & \texttt{cube}, \texttt{form}, \texttt{trans} \\
Audio & 6 & \texttt{mic}, \texttt{melody}, \texttt{overtone} \\
Control & 7 & \texttt{def}, \texttt{if}, \texttt{repeat}, \texttt{once}, \texttt{later} \\
Text & 4 & \texttt{text}, \texttt{font} \\
Input & 3 & \texttt{pen}, \texttt{touch} \\
Composition & 4 & \texttt{embed}, \texttt{layer}, \texttt{fade} \\
Data & 4 & \texttt{let}, \texttt{list}, \texttt{get}, \texttt{set} \\
\midrule
\textbf{Total} & \textbf{118} & \\
\bottomrule
\end{tabularx}
\caption{KidLisp built-in functions by category.}
\label{tab:functions}
\end{table}

\subsection{Syntax Examples}

A complete \kidlisp{} program that draws a rotating red circle:

\begin{lstlisting}
(wipe black)
(spin (* frame 0.01))
(ink red)
(circle (/ width 2) (/ height 2) 50)
\end{lstlisting}

An animated gradient background with embedded audio:

\begin{lstlisting}
(wipe (rainbow (* frame 0.5)))
(melody "c4 e4 g4 c5")
(ink white 128)
(text "hello" (/ width 2) (/ height 2))
\end{lstlisting}

A program that embeds another program as a layer:

\begin{lstlisting}
(wipe black)
(layer $cow 0.5)
(spin 0.5)
(zoom 1.01)
\end{lstlisting}

Here \texttt{\$cow} is a reference to another \kidlisp{} program by its three-character code. The \texttt{layer} function composites it with 50\% opacity.

\subsection{What KidLisp Omits}

The language deliberately excludes:

\begin{itemize}
  \item User-defined functions (beyond \texttt{def} for named values)
  \item Recursion
  \item String manipulation
  \item File I/O and networking
  \item Mutable data structures
  \item Exception handling
\end{itemize}

These omissions are not limitations---they are the design. By removing abstraction mechanisms, \kidlisp{} keeps programs flat and legible. A reader can understand any program by reading it top-to-bottom, without tracing call stacks. Composition happens through embedding, not through abstraction.

% ============ 3. ARCHITECTURE ============

\section{Evaluator Architecture}
\label{sec:architecture}

The evaluator (\texttt{kidlisp.mjs}, approximately 50KB) is a single JavaScript ES module implementing a tree-walking interpreter.

\subsection{Evaluation Model}

\begin{enumerate}
  \item \textbf{Parse}: S-expressions are tokenized and parsed into a nested array AST.
  \item \textbf{Walk}: The evaluator traverses the tree, dispatching on the first symbol of each list.
  \item \textbf{Execute}: Built-in functions map directly to Canvas 2D operations, Web Audio API calls, and mathematical computations.
\end{enumerate}

The evaluator runs once per animation frame (typically 60 FPS). The entire program is re-evaluated each frame, with \texttt{frame} and \texttt{time} bindings updated automatically. This model---borrowed from immediate-mode graphics---eliminates the need for explicit animation loops or state management.

\subsection{Deterministic Rendering}

Given the same random seed and frame count, a \kidlisp{} program always produces identical output. The \texttt{random} function uses a seeded PRNG initialized from the program's short code. This enables reproducible generative art: the same code at the same frame always looks the same.

\subsection{Embedding and Composition}

The \texttt{embed} and \texttt{layer} functions load other \kidlisp{} programs by their short codes. The evaluator maintains a cache of parsed programs and renders them into offscreen canvases, compositing the results with configurable opacity and blend modes.

This creates a composition graph: programs reference other programs, forming chains of visual layers. In practice, users build complex visuals by combining simple programs rather than writing complex single programs.

\subsection{Performance}

The evaluator processes typical programs in under 1ms per frame. Transformation functions (\texttt{blur}, \texttt{zoom}, \texttt{spin}) operate on pixel buffers and are the primary performance cost. Embedding adds the cost of evaluating the referenced program, but caching mitigates this for static programs.

% ============ 4. SOCIAL DISTRIBUTION ============

\section{Social Distribution}
\label{sec:social}

Every \kidlisp{} program, upon creation, is assigned a unique three-character alphanumeric code (e.g., \texttt{cow}, \texttt{27z}, \texttt{gfl}) and stored in MongoDB via a REST API.

\subsection{Sharing Mechanisms}

Programs are shareable through multiple channels:

\begin{itemize}
  \item \textbf{URL}: \texttt{aesthetic.computer/\$cow} renders the program in a browser
  \item \textbf{Code}: Users type \texttt{\$cow} in the Aesthetic Computer prompt to view any program
  \item \textbf{Embedding}: \texttt{(\$cow)} inside another program includes it as a visual layer
  \item \textbf{QR code}: \texttt{share \$cow} generates a scannable code
\end{itemize}

The three-character code is memorable and speakable. Users share programs verbally (``check out dollar cow''), in chat, and by QR code at physical events. This is intentional: the distribution mechanism is designed for social contexts, not package managers.

\subsection{Remixing}

Any program's source is viewable. Users routinely copy a program's code, modify it, and save a new version. This creates informal lineages of remixed programs---a pattern well-documented in Scratch~\citep{resnick2009scratch} and Shadertoy communities.

% ============ 5. EVALUATION ============

\section{Evaluation}
\label{sec:evaluation}

\kidlisp{} has been available on the Aesthetic Computer platform since 2024. As of March 2026:

\begin{table}[h]
\small
\centering
\begin{tabular}{lr}
\toprule
\textbf{Metric} & \textbf{Value} \\
\midrule
Total programs created & 16,174 \\
Unique authors & 59 \\
Programs per author (median) & 12 \\
Programs per author (top quartile) & 180+ \\
Embedding depth (max observed) & 7 layers \\
Platform registered users & 2,798 \\
\bottomrule
\end{tabular}
\caption{KidLisp adoption metrics (March 2026).}
\label{tab:adoption}
\end{table}

\subsection{Observations}

The 16,174 programs represent a naturalistic corpus of creative coding practices produced by a non-expert population. Several patterns emerge:

\textbf{Composition over abstraction.} Users embed rather than abstract. The most-referenced programs are simple visual primitives (gradients, shapes, patterns) that serve as building blocks for more complex compositions.

\textbf{Exploration over engineering.} Programs are typically short (under 10 lines). Users iterate by saving new programs rather than editing existing ones, creating divergent exploration paths.

\textbf{Social learning.} The most prolific authors began by viewing and remixing existing programs. The three-character code system makes discovery accidental---users encounter programs by trying random codes.

\subsection{Limitations}

\kidlisp{} is not suitable for general-purpose programming, complex simulations, or applications requiring state persistence across sessions. Its 118-function vocabulary is sufficient for generative art but would be constraining for other domains. The lack of user-defined functions limits program complexity---a deliberate choice, but one that some advanced users find frustrating.

% ============ 6. RELATED WORK ============

\section{Related Work}

\kidlisp{} sits at the intersection of creative coding environments, educational programming languages, and live coding tools.

\textbf{Creative coding.} Processing~\citep{reas2007processing} and p5.js~\citep{mccarthy2015p5js} are the dominant tools, using Java and JavaScript respectively. OpenFrameworks uses C++. All require managing a project structure and understanding a general-purpose language. \kidlisp{} trades generality for immediacy.

\textbf{Live coding.} Hydra~\citep{hydra2019} provides a JavaScript-based live visual coding environment. TidalCycles uses Haskell for music. Sonic Pi~\citep{aaron2016sonic} uses Ruby for music education. Fluxus~\citep{fluxus2007} used a Scheme-like language for 3D visuals. \kidlisp{} is closest to Fluxus in spirit but targets 2D generative art in the browser.

\textbf{Educational languages.} Scratch~\citep{resnick2009scratch} demonstrated that social sharing drives engagement in educational programming. Logo introduced turtle graphics as a minimal drawing primitive. \kidlisp{} inherits Logo's philosophy of immediate visual feedback but uses a Lisp syntax and targets generative art rather than geometry.

\textbf{Lisp for art.} Racket's \texttt{2htdp/image} library~\citep{racket2018} provides functional image construction for CS education. Context Free Art uses a grammar-based approach. \kidlisp{} differs by targeting real-time animation and embedding a social distribution mechanism directly in the language.

% ============ 7. CONCLUSION ============

\section{Conclusion}

\kidlisp{} demonstrates that a minimal domain-specific language---118 functions, no user-defined abstractions, browser-native---can support a productive community of generative art practitioners. The key design insight is that social distribution (three-character codes, URL sharing, embedding) is not an add-on but a core language feature that shapes how users learn, create, and build on each other's work.

The corpus of 16,174 programs created by 59 authors provides evidence that non-expert users can engage productively with a Lisp-family language when the domain is constrained and the feedback is immediate. We suggest that this design point---trading computational generality for social accessibility---deserves further exploration in the language design community.

\kidlisp{} is open source under the ISC license. The evaluator, documentation, and tools are available at \url{https://github.com/digitpain/aesthetic.computer}. Programs can be browsed and created at \url{https://kidlisp.com}.

\vspace{0.5em}
\noindent\textbf{ORCID:} \href{https://orcid.org/0009-0007-4460-4913}{0009-0007-4460-4913}

% ============ REFERENCES ============

\bibliographystyle{plainnat}
\bibliography{references}

\end{document}
